\section{Consequences and Decidability of Convertibility}

In the previous section, we have established the fundamental theorems, and using the
escape lemma, we see that all syntactic equivalent types and terms are also
semantically related by the Kripke logical relations. %
In this section, we will instantiate the generic equivalence so that we can obtain
desired properties that are difficult to prove syntactically. %

\subsection{First Instantiation: Syntactic Equivalence}

First we instantiate the generic equivalence with syntactic equivalences of types and
terms. %
The laws are all easily instantiated by the corresponding equivalence rules. %
We also see that the derived equivalences between local contexts and local
substitutions are equivalent to the corresponding syntactic equivalences. %
From the fundamental theorems and the escape lemma, we are able to derive the
following lemmas.

\begin{lemma}[Injectivity of Type Constructors] $ $
  \begin{itemize}
  \item If $\lttypeq i{\PI{l}{l'} x S T}{\PI{l}{l'} x{S'}{T'}}{l \sqcup l'}$ and
    $i \in \{\proglevel, \metalevel\}$, then $\lttypeq i{S}{S'}l$ and
    $\lttypeq[\Psi][\Gamma, x : S \at l]i{T}{T'}{l'}$.
  \item If
    $\lttyequiv i{\PI{l}{l'} x s t}{\PI{l}{l'} x{s'}{t'}}{\Ty{l \sqcup l'}}{\su{(l
        \sqcup l')}}$ and $i \in \{\proglevel, \metalevel\}$, then \break
    $\lttyequiv i{s}{s'}{\Ty l}{1 + l}$ and
    $\lttyequiv[\Psi][\Gamma, x : \Elt l s \at l]i{t}{t'}{\Ty{l'}}{1 + l'}$.
  \item If $\lttypeq \metalevel{\CTyp[\Delta]l}{\CTyp[\Delta']l}{0}$, then
    $\lpequiv \proglevel \Delta{\Delta'}$.
  \item If $\lttypeq \metalevel{\CTrm[\Delta] T l}{\CTrm[\Delta'] T l}{0}$, then
    $\lpequiv \proglevel \Delta{\Delta'}$ and $\lttypeq[\Psi][\Delta] \proglevel{T}{T'}l$.
  \item If $\lttypeq \metalevel{\CPI g l T}{\CPI g l{T'}}l$, then
    $\lttypeq[\Psi, g : \Ctx] \metalevel{T}{T'}l$.
  \item If
    $\lttypeq \metalevel{\TPI{U}[\Delta]{l}{l'}{T}}{\TPI{U}[\Delta']{l}{l'}{T'}}{l'}$, then $\lpequiv \proglevel \Delta{\Delta'}$ and
    $\lttypeq[\Psi, U : \DTyp[\Delta] \proglevel l] \metalevel{T}{T'}{l'}$.
  \item If $\lttypeq \metalevel{\UPI \ell l T}{\UPI \ell l{T'}}\omega$, then
    $\lttypeq[\Psi][\Gamma][L, \vect \ell] \metalevel{T}{T'}{l}$.
  \end{itemize}
\end{lemma}
\begin{proof}
  By fundamental theorems, we pass in all identity substitutions and then we can
  extract this fact from the logical relations of types. 
\end{proof}

We follow \citet{hu_jang_pientka_2023} and give us the consistency proof of \delamlang.
\begin{lemma}[Consistency]
  There is no term $t$ that satisfies this typing judgment:
  \[
    \lttyping[\cdot][\cdot][\ell] \metalevel t{\PI{1 + \ell}{\ell}x{\Ty \ell}x}{1 + \ell}
  \]
\end{lemma}
That is, there is not a generic way to construct a term of an arbitrary type.
\begin{proof}
  The lemma effectively asks to reject the following derivation after applying the
  fundamental theorems:
  \[
    \lttyping[\cdot][x : \Ty \ell \at{1 + \ell}][\ell] \metalevel t x \ell
  \]
  Now by the logical relations of the neutral type $x$, we know that $t$ must also be
  neutral, so we now move on to rejecting 
  \[
    \lttyping[\cdot][x : \Ty \ell \at{1 + \ell}][\ell] \metalevel \varlevel x \ell
  \]
  Then we do induction on $\varlevel$. %
  It is impossible to do any operation on $\varlevel$ but to refer to $x$ ultimately, but then
  $x$ has type $\Ty \ell$ which cannot be equivalent to $x$, as they both have reached
  normal forms and it is not possible to relate any two distinguished normal forms in
  any case by the logical relations. %
\end{proof}

\subsection{Second Instantiation: Convertibility Checking}

In this section, we perform our second instantiation by specifying the generic
equivalence to be the convertibility checking judgments. %
This second instantiation is not very immediate for some laws for the generic
equivalence, so we must do some verification, using the results from the first
instantiation. %

\begin{itemize}
\item $\lttypgneeq i V{V'} l$ is $\tconvtypne i V{V'}l$.
\item $\lttypgeq i {T}{T'} l$ is $\tconvtyp i T{T'} l$.
\item $\lttrmgneeq i \nu{\nu'} T l$ is $\tconvtrmne i \nu{\nu'}{T'} l$ and
  $\lttypeq i T{T'} l$.
\item $\lttrmgeq i t{t'} T l$ is $\tconvtrm i t{t'}T l$.
\end{itemize}

Most laws are quite straightforward, except the conversion laws and the PER laws. %
Let us consider the conversion laws first. %
We must show that the convertibility checking algorithm is invariant under contexts
and types. %
Intuitively, this should be true, but it is not very easy to prove, until we obtain
the injectivity lemma after the first instantiation.

\begin{lemma}[Conversion] $ $
  \begin{itemize}
  \item if $\tconvtyp i T{T'} l$, $\gequiv\Phi\Psi$ and
    $\lpequiv[\Phi]{i}\Delta\Gamma$, then $\tconvtyp[\Phi][\Delta] i T{T'} l$.
  \item if $\tconvtypnf i W{W'} l$, $\gequiv\Phi\Psi$ and
    $\lpequiv[\Phi]{i}\Delta\Gamma$, then $\tconvtypnf[\Phi][\Delta] i W{W'} l$.
  \item if $\tconvtypne i V{V'}l$, $\gequiv\Phi\Psi$ and
    $\lpequiv[\Phi]{i}\Delta\Gamma$, then $\tconvtypne[\Phi][\Delta] i V{V'}l$.
  \item if $\tconvctx i \Gamma \Delta$ and $\gequiv\Phi\Psi$, then $\tconvctx[\Phi] i \Gamma \Delta$.
  \item if $\tconvtrm i t{t'}T l$, $\gequiv\Phi\Psi$, $\lpequiv[\Phi]{i}\Delta\Gamma$
    and $\lttypeq[\Phi][\Delta] i{T'} T l$, then
    $\tconvtrm[\Phi][\Delta] i t{t'}{T'} l$.
  \item if $\tconvtrmnf i w{w'}W l$, $\gequiv\Phi\Psi$, $\lpequiv[\Phi]{i}\Delta\Gamma$
    and $\lttypeq[\Phi][\Delta] i{W'}W l$, then $\lttyequiv[\Phi][\Delta] i w{w'}{W'} l$;
  \item if $\tconvtrmne i \nu{\nu'}T l$, $\gequiv\Phi\Psi$ and
    $\lpequiv[\Phi]{i}\Delta\Gamma$, then $\tconvtrmne[\Phi][\Delta] i \nu{\nu'}{T'} l$
    and $\lttypeq i{T'}{T}l$ for some $T'$. 
  \item if $\tconvtrmnee i \nu{\nu'}W l$, $\gequiv\Phi\Psi$ and
    $\lpequiv[\Phi]{i}\Delta\Gamma$, then $\lttyequiv[\Phi][\Delta] i \nu{\nu'}{W'} l$
    and $\lttypeq i{W'}{W}l$ for some $W'$.
  \item if $\tconvsub i{\delta}{\delta'}\Delta$, $\gequiv\Phi\Psi$,
    $\lpequiv[\Phi]{i}{\Gamma'}\Gamma$ and $\lpequiv[\Phi]{i}{\Delta'}\Delta$, then
    $\tconvsub[\Phi][\Gamma'] i{\delta}{\delta'}{\Delta'}$.
  \end{itemize}
\end{lemma}
\begin{proof}
  We do induction on all derivations. %
  Most cases are immediate. %
  When we get under binders, we need to extend the context equivalences, in which case
  we should use the soundness lemma to obtain the well-formedness of the type that
  needs to be extended to the contexts. %

  We consider a few cases.
  \begin{itemize}[label=Case]
  \item
    \begin{mathpar}
      \inferrule
      { T \reds W \\ \ttrmreds i t w T l \\ \ttrmreds i
        {t'}{w'}T l \\ \tconvtrmnf i w{w'}W l}
      {\tconvtrm i t{t'}T l}
    \end{mathpar}
    The most important thing to notice is that by $\lttypeq[\Phi][\Delta] i{T'} T l$
    and $T \reds W$, we know that $T' \reds W'$ for some $W'$. Therefore we can prove
    this case using IH. %
    Other premises are satisfied by syntactic context equivalence lemmas. %
    
  \item
    \begin{mathpar}
      \inferrule
      {\lttypwf i S l \\ \lttyping i w {\PI{l}{l'}x{S}{T}}{l \sqcup l'} \\ \lttyping i {w'}{\PI{l}{l'}x{S}{T}}{l \sqcup l'} \\ \tconvtrm[\Psi][\Gamma, x : S \at{l}] i {\APP{w}{l}{l'}x{S}{T}{x}}{\APP{w'}{l}{l'}x{S}{T}{x}}{T}{l'}}
      {\tconvtrmnf i{w}{w'}{\PI{l}{l'}x{S}{T}}{l \sqcup l'}}
    \end{mathpar}
    In this case, we know $\lttypeq[\Phi][\Delta] i{W'}{\PI{l}{l'}x{S}{T}}{l \sqcup
      l'}$. %
    By the fundamental theorems, we know that $W'$ can only reduce to some $\Pi$
    type. %
    Since $W'$ is already normal, it is only possible for $W'$ to be some $\Pi$
    type. %
    Say $W' = {\PI{l}{l'}x{S'}{T'}}$. %
    Then by injectivity, we know $S \approx S'$ and $T \approx T'$. %
    We obtain our goal by extending the local contexts. 
    
  \item
    \begin{mathpar}
      \inferrule
      {\tconvtrmnee i \nu{\nu'}{W} l}
      {\tconvtrmnf i \nu{\nu'}V l}
    \end{mathpar}
    In this case, the type equivalence is irrelevant. %
    When hitting neutral types, we simply ignore the type $W$ inferred by
    $\tconvtrmnee i \nu{\nu'}{W} l$.

  \item
    \begin{mathpar}
      \inferrule
      {\tconvtrmne i \nu{\nu'}{T}{l} \\ T \reds W}
      {\tconvtrmnee i \nu{\nu'}W l}
    \end{mathpar}
    In this case, we simply apply IH. Then we know
    \[
      T \reds W \tand T' \reds W'
    \]
    Our goal is to show $W \approx W'$. %
    But this is immediate from the determinacy lemma of multi-step reduction and the
    fundamental theorems. %

  \item
    \begin{mathpar}
      \inferrule
      {\lpjudge i \Gamma \\ x : T \at l \in \Gamma}
      {\tconvtrmne i x{x}T l}
    \end{mathpar}
    The goal is given by the equivalence between $\Gamma$ and $\Delta$. 
    
  \item
    \begin{mathpar}
      \inferrule
      {\tyequiv[L]{l_1}{l_3}\Level  \\ \tyequiv[L]{l_2}{l_4}\Level \\ \tconvtyp i S{S'}{l_1}
        \\
        \tconvtyp[\Psi][\Gamma, x : S \at{l_1}] i{T}{T'}{l_2} \\
        \tconvtrmnee i \nu{\nu'}{\PI {l_1}{l_2} x{S''}{T''}}{l_1 \sqcup l_2} \\ \tconvtrm i s{s'} S{l_1}}
      {\tconvtrmne i{\APP \nu {l_1}{l_2} x S T s}{\APP{\nu'}{l_3}{l_4} x {S'}{T'}{s'}}{T[s/x]}{l_2}}
    \end{mathpar}
    By IH, we know from
    $\tconvtrmnee i \nu{\nu'}{\PI {l_1}{l_2} x{S''}{T''}}{l_1 \sqcup l_2}$ that there
    must be $W'$, so that
    \[
      \lttypeq i{W'}{\PI {l_1}{l_2} x{S''}{T''}}{l_1 \sqcup l_2}
    \]
    By the fundamental theorems, we know $W'$ must be some $\Pi$ types. %
    The return type is fixed so we do not have to do anything. %
  \end{itemize}
\end{proof}

From the conversion lemma, we see that both $\tconvtrmne i \nu{\nu'}T l$ and
$\tconvtrmnee i \nu{\nu'}W l$ return some equivalent types because they are in fact
inference steps. %
Therefore, it makes sense when we instantiate $\lttrmgneeq i \nu{\nu'} T l$, we hide a
syntactic equivalence judgment inside. %

For the PER laws, we see that the difficulties primarily come from the convertibility
checking of neutrals, because again they are inference steps. %
Moreover, due to their left-biased design, in general it is not true that the inferred
types can be replaced by their equivalence. %
However, since we are hiding an equivalence judgment in $\lttrmgneeq i \nu{\nu'} T l$
during instantiation, following the same principle as the conversion lemma, we are
able to erase the effect of the left bias and establish the PER laws. %

The instantiation immediately gives us the convertibility lemma once we apply the
escape lemma:
\begin{theorem}[Convertibility] $ $
  \begin{itemize}
  \item If $\lttypwf i T l$, then $\tconvtyp i T T l$. 
  \item If $\lttypeq i T{T'} l$, then $\tconvtyp i T{T'} l$. 
  \item If $\lttyping i t T l$, then $\tconvtrm i t t T l$.
  \item If $\lttyequiv i t{t'} T l$, then $\tconvtrm i t{t'} T l$.
  \end{itemize}
\end{theorem}
In particular, we see that the convertibility checking algorithm is both sound and
complete w.r.t. the syntactic equivalence judgments. %
The decidability of type checking requires the following lemma by attempting to relate
two reflexively convertible types or terms. 

\begin{lemma}[Decidability] $ $
  \begin{itemize}
  \item if $\tconvtyp[\Phi][\Delta] i T{T} l$, $\tconvtyp i {T'}{T'} l$,
    $\gequiv\Phi\Psi$ and $\lpequiv[\Phi]{i}\Delta\Gamma$, then whether
    $\tconvtyp[\Phi][\Delta] i T{T'} l$ is decidable.
  \item if $\tconvtypnf[\Phi][\Delta] i W{W} l$, $\tconvtypnf i{W'}{W'} l$,
    $\gequiv\Phi\Psi$ and $\lpequiv[\Phi]{i}\Delta\Gamma$, then whether
    $\tconvtypnf[\Phi][\Delta] i W{W'} l$ is decidable.
  \item if $\tconvtypne[\Phi][\Delta] i V{V}l$, $\tconvtypne i {V'}{V'}l$,
    $\gequiv\Phi\Psi$ and $\lpequiv[\Phi]{i}\Delta\Gamma$, then whether
    $\tconvtypne[\Phi][\Delta] i V{V'}l$ is decidable. %
  \item if $\tconvctx[\Phi] i \Gamma \Gamma$, $\tconvctx i \Delta \Delta$ and
    $\gequiv\Phi\Psi$, then whether $\tconvctx[\Phi] i \Gamma \Delta$ is decidable. 
  \item if $\tconvtrm[\Phi][\Delta] i t{t}T l$, $\tconvtrm i{t'}{t'}{T} l$,
    $\gequiv\Phi\Psi$ and $\lpequiv[\Phi]{i}\Delta\Gamma$, then whether
    $\tconvtrm[\Phi][\Delta] i t{t'}{T} l$ is decidable.
  \item if $\tconvtrmnf[\Phi][\Delta] i w{w}W l$, $\tconvtrmnf i {w'}{w'}{W} l$,
    $\gequiv\Phi\Psi$ and $\lpequiv[\Phi]{i}\Delta\Gamma$, then whether
    $\lttyequiv[\Phi][\Delta] i w{w'}{W} l$ is decidable.
  \item if $\tconvtrmne[\Phi][\Delta] i \nu{\nu}T l$,
    $\tconvtrmne i {\nu'}{\nu'}{T'} l$, $\gequiv\Phi\Psi$ and
    $\lpequiv[\Phi]{i}\Delta\Gamma$, then whether
    $\tconvtrmne[\Phi][\Delta] i \nu{\nu'}{T''} l$ for some $T''$ is decidable.
  \item if $\tconvtrmnee[\Phi][\Delta] i \nu{\nu}W l$,
    $\tconvtrmnee i {\nu'}{\nu'}{W'} l$, $\gequiv\Phi\Psi$ and
    $\lpequiv[\Phi]{i}\Delta\Gamma$, then whether
    $\lttyequiv[\Phi][\Delta] i \nu{\nu'}{W''} l$ is decidable.
  \item if $\tconvsub[\Phi][\Gamma'] i{\delta}{\delta}{\Delta}$,
    $\tconvsub i{\delta'}{\delta'}{\Delta}$, $\gequiv\Phi\Psi$ and
    $\lpequiv[\Phi]{i}{\Gamma'}\Gamma$, then whether
    $\tconvsub[\Phi][\Gamma'] i{\delta}{\delta'}{\Delta}$ is decidable. %
  \end{itemize}
\end{lemma}
\begin{proof}
  We do a mutual induction on the first derivations and then invert the second ones. %
  We can reject most cases when they have different root derivations. %
  We consider a few cases.
  \begin{itemize}[label=Case]
  \item
    \begin{mathpar}
      \inferrule
      {\ttypreds[\Phi][\Delta] i T W l \\\\ \tconvtypnf i W{W} l}
      {\tconvtyp[\Phi][\Delta] i T{T} l}

      \inferrule
      {\ttypreds i{T'}{W'} l \\\\ \tconvtypnf i{W'}{W'} l}
      {\tconvtyp i T{T'} l}
    \end{mathpar}
    In this case, we apply IH to decide whether $W$ and $W'$ are convertible. %

  \item
    \begin{mathpar}
      \inferrule
      {\tconvtyp[\Phi][\Delta] i S{S} l \\ \tconvtyp[\Phi][\Delta, x : S \at l] i T{T}{l'}}
      {\tconvtypnf[\Phi][\Delta] i{\PI l{l'}x{S}{T}}{\PI l{l'}x{S}{T}}{l \sqcup l'}}

      \inferrule
      {\tconvtyp i{S'}{S'} l \\ \tconvtyp[\Psi][\Gamma, x : S' \at l] i{T'}{T'}{l'}}
      {\tconvtypnf i{\PI l{l'}x{S'}{T'}}{\PI l{l'}x{S'}{T'}}{l \sqcup l'}}
    \end{mathpar}
    We can decide whether $S$ and $S'$ are convertible by IH. %
    When we decide $T$ and $T'$, we must extend the equivalent local contexts. %
    
  \item
    \begin{mathpar}
      \inferrule
      {T \reds W \\ \ttrmreds[\Phi][\Delta] i t w T l \\\\  \tconvtrmnf[\Phi][\Delta] i w{w}W l}
      {\tconvtrm[\Phi][\Delta] i t{t}T l}

      \inferrule
      {T' \reds W' \\ \ttrmreds i{t'}{w'}{T'} l \\\\  \tconvtrmnf i {w'}{w'}{W'} l}
      {\tconvtrm i {t'}{t'}{T} l}
    \end{mathpar}
    By fundamental theorems, we know $W \approx W'$. %
    Then by IH, we can decide whether $w$ and $w'$ are convertible. %
    
  \item
    \begin{mathpar}
      \inferrule
      {\lttypwf[\Phi][\Delta] i S l \\ \lttyping[\Phi][\Delta] i w
        {\PI{l}{l'}x{S}{T}}{l \sqcup l'} \\
        \tconvtrm[\Phi][\Delta, x : S \at{l}] i {\APP{w}{l}{l'}x{S}{T}{x}}{\APP{w}{l}{l'}x{S}{T}{x}}{T}{l'}}
      {\tconvtrmnf[\Phi][\Delta] i{w}{w}{\PI{l}{l'}x{S}{T}}{l \sqcup l'}}

      \inferrule
      {\lttypwf i{S} l \\ \lttyping i{w'}{\PI{l}{l'}x{S}{T}}{l \sqcup l'} \\
        \tconvtrm[\Psi][\Gamma, x : S \at{l}] i {\APP{w'}{l}{l'}x{S}{T}{x}}{\APP{w'}{l}{l'}x{S}{T}{x}}{T}{l'}}
      {\tconvtrmnf i{w'}{w'}{\PI{l}{l'}x{S}{T}}{l \sqcup l'}}
    \end{mathpar}
    Again, it is quite immediate by IH. %
    
  \item
    \begin{mathpar}
      \inferrule
      {\tconvtrmnee[\Phi][\Delta] i \nu{\nu}{W} l}
      {\tconvtrmnf[\Phi][\Delta] i \nu{\nu}V l}

      \inferrule
      {\tconvtrmnee i {\nu'}{\nu'}{W'} l}
      {\tconvtrmnf i {\nu'}{\nu'}{V} l}
    \end{mathpar}
    By IH, we know $\tconvtrmnee[\Phi][\Delta] i \nu{\nu'}{W''} l$ for some $W''$. 

  \item
    \begin{mathpar}
      \inferrule
      {\tconvtrmne[\Phi][\Delta] i \nu{\nu}{T}{l} \\ T \reds W}
      {\tconvtrmnee[\Phi][\Delta] i \nu{\nu}W l}

      \inferrule
      {\tconvtrmne i {\nu'}{\nu'}{T'}{l} \\ T' \reds W'}
      {\tconvtrmnee i {\nu'}{\nu'}{W'} l}
    \end{mathpar}
    By IH, we have $\tconvtrmne i {\nu'}{\nu'}{T''}{l}$ for some $T''$. %
    $T''$ will reduce to some normal form by the fundamental theorems. %
  \end{itemize}
\end{proof}

\begin{theorem}[Decidability of Convertibility] $ $
  \begin{itemize}
  \item If $\lttypwf i T l$ and $\lttypwf i{T'} l$, then whether $\tconvtyp i T{T'} l$
    is decidable.
  \item If $\lttyping i t T l$ and $\lttyping i{t'} T l$, then whether
    $\tconvtrm i t{t'}T l$ is decidable.
  \end{itemize}
\end{theorem}
\begin{proof}
  First we use the fundamental theorems from the second instantiation to show that
  both types (or terms, resp.) are reflexively convertible. %
  Then we use the decidability lemma above. %
\end{proof}

At this point, we have justified the decidability of convertibility checking of
\delamlang and therefore conclude our investigations. %
